\documentclass[a4paper,12pt]{article}
\usepackage{amsmath}
\usepackage[margin=1in]{geometry}
\begin{document}

\title {Circular motion on an inclined plane with friction}
\maketitle
		

Consider a car turning on a bank with an incline. The car's mass is 1000 kg, and the radius of curvature is 50 m. Let the coefficient of static friction be $\mu_s = 0.3$ and the angle on inclination be $\theta = 20^\circ $.



\begin{enumerate} 
	\item{In standard horizontal and vertical coordinates, write out an equation for the vertical equalibrium and another for horizontal equilibrium. The vertical equation should be equal to the mass times the gravitational acceleration. Likewise the horizontal equation should be equal to the mass times the centripetal acceleration.}
	\item{Note that the direction of static friction is going to depend on if the car is sliding outwards or falling inwards. Write an equation for the maximum velocity of the car without sliding out. To do this, isolate for mass in your two equations, and set them equal. The normal forces should cancel out.}
	\item{Now find an equation for the minimum velocity to maintain stability without falling inward. Why is the friction used here static and not kinetic?}
	\item{Now set the coefficient of static friction to be zero. What is your new equation for velocity? It should have tan in it. This is the equation you need to answer the bicycle question from last week.}
	\item{Solve for the min and max velocities with the numbers given in the start of the question.}

\end{enumerate}


\end{document}
